\documentclass[12pt, fleqn]{article}

\usepackage[utf8]{inputenc}
\usepackage[T2A]{fontenc}
\usepackage{amssymb, amsmath, mathrsfs, amsthm}
\usepackage[russian]{babel}
\usepackage{graphicx}
\usepackage[footnotesize]{caption2}
\usepackage{indentfirst}
\usepackage[hidelinks]{hyperref}
\usepackage{multirow}
%\usepackage[ruled,section]{algorithm}
%\usepackage[noend]{algorithmic}
%\usepackage[all]{xy}

% Параметры страницы
\textheight=24cm
\textwidth=16cm
\oddsidemargin=5mm
\evensidemargin=-5mm
\marginparwidth=36pt
\topmargin=-1cm
\footnotesep=3ex
%\flushbottom
\raggedbottom
\tolerance 3000
% подавить эффект "висячих стpок"
\clubpenalty=10000
\widowpenalty=10000
\renewcommand{\baselinestretch}{1.1}
\renewcommand{\baselinestretch}{1.5} %для печати с большим интервалом

% Дополнительные команды для личных обозначений
\newcommand{\expectation}{\mathop{\mathbb{E}}}
\newcommand{\norm}[1]{\left\lVert#1\right\rVert}
\newcommand{\loss}{\mathop{\mathcal{L}}}
\newcommand{\mse}{\mathop{MSE}}
\newcommand{\scalarproduct}[1]{\langle #1 \rangle}

\newcommand{\predictionfunction}{\hat{f}}
\newcommand{\ensemblefunction}{a}
\newcommand{\optimizationmethodfunction}{\tilde{f}}
\newcommand{\distinguishparameter}{\alpha}
\newcommand{\objects}{X}
\newcommand{\results}{Y}
\newcommand{\predictedobjects}{\widehat{\objects}}

\newcommand{\numberobjects}{N}
\newcommand{\numberpredictionfunctions}{M}

\newcommand{\for}[3]{\sum\limits_{#1 = #2}^{#3}}  % Usage: \for{index}{begin}{end}
\newcommand{\forn}[2]{\for{#1}{1}{#2}}  % Usage: \forn{index}{end}

\newcommand{\many}[3]{#1 1 #2, #1 2 #2, \dots, #1 #3 #2}  % Usage: \many{prefix}{suffix}{end}

\newcommand{\reference}[1]{(\hyperref[#1]{\ref{#1}})}

\newcommand{\ensemblefunctionfull}{\ensemblefunction(\many{\predictionfunction_}{(x)}{\numberpredictionfunctions})}

\begin{document}

\begin{titlepage}
\begin{center}
    Московский государственный университет имени М. В. Ломоносова

    \bigskip
    \includegraphics[width=50mm]{msu.eps}

    \bigskip
    Факультет Вычислительной Математики и Кибернетики\\
    Кафедра Математических Методов Прогнозирования\\[10mm]

    \textsf{\large\bfseries
        <<Сравнение скорости вычисления собственных значений положительно определённых матриц>>
    }\\[20mm]

    \begin{flushright}
        \parbox{0.5\textwidth}{
            Выполнил:\\
            студент 1 курса магистратуры 517 группы\\
            \emph{Королев Николай Сергеевич}\\[5mm]
            Преподаватель:\\
            канд. техн. наук, доцент\\
            \emph{Русол Андрей Владимирович}
        }
    \end{flushright}

    \vspace{\fill}
    Москва, 2021
\end{center}
\end{titlepage}

\newpage
\renewcommand{\contentsname}{Содержание}
\tableofcontents

\newpage
\begin{abstract}
    В данной работе рассматривается новый метод повышения качества обучения моделей, а также увеличения их обобщающей способности, основанный на ансамбле промежуточных решений. В ходе вычислительного эксперимента, была показана возможная применимость данного метода для улучшения качества нейронных сетей, решающих задачу классификации на реальных данных.
\end{abstract}

\newpage
\section{Введение}

В настоящее время придумано огромное количество алгоритмов машинного обучения. Данные алгоритмы решают различные задачи и применяются во многих сферах человеческой деятельности. Одной из наиболее распространённых способов применения машинного обучения является решение задачи предсказания, в которой необходимо предсказать некоторый параметр (отклик) по ряду других параметров. В таких задачах крайне важна обобщающая способность используемого алгоритма обучения. В случае недостаточной обобщающей способности, итоговая модель не сможет достаточно точно предсказывать необходимый параметр (отклик), поэтому очень остро стоит проблема повышения обобщающей способности модели.

В данной работе рассматривается новый метод повышения качества обучения моделей, а также увеличения их обобщающей способности, основанный на ансамбле промежуточных решений. Кроме того, проводятся вычислительные эксперименты, показывающие эффективность работы данного метода в задаче классификации на данных систолического артериального давления.

\subsection{Постановка задачи}

$\many{x_}{}{\numberobjects}$~--~точки в некотором векторном пространстве; $\many{y_}{}{\numberobjects}$~--~ значения, соответствующие этим точкам. При этом $y_i = y(x_i) = f(x_i) + \varepsilon$, где $\varepsilon$~--~случайная величина с нулевым математическим ожиданием и дисперсией $\sigma^2$. Пусть также задана функция $\predictionfunction_1(x)$, которая приближает функцию $f(x)$. Стоит задача нахождения $\numberpredictionfunctions - 1$ функции $\predictionfunction_2(x), \predictionfunction_3(x), \dots, \predictionfunction_\numberpredictionfunctions(x)$, а также функции-ансамбля $\predictionfunction(x) = \ensemblefunctionfull$, приближающей функцию $f(x)$ лучше чем любая из функций $\predictionfunction_i(x), \; i = \many{}{}{\numberpredictionfunctions}$.

\subsection{Теоретическая часть}
Из~\cite{BiasVarianceDecompositionZeroOneSquaredLoss} известно разложение математического ожидания квадратичной ошибки на смещение и дисперсию:
\begin{equation*}
\expectation \left[ \left( \predictionfunction(x) - y \right)^2 \right] =
\left(
	\expectation \predictionfunction(x) - f(x)
\right)^2  +
\left(
	\expectation \predictionfunction^2(x) - \left(\expectation \predictionfunction(x)\right)^2
\right) + \sigma^2
\end{equation*}

Для уменьшения как смещения, так и дисперсии $\predictionfunction(x)$ используются ансамбли:
\begin{equation*}
\predictionfunction(x) = \ensemblefunctionfull
\end{equation*}

В~\cite{OptimalConvexCorrectingProcedures} было доказано, что в случае, когда функция $\ensemblefunctionfull$ является выпуклой комбинацией своих аргументов:
\begin{equation*}
\ensemblefunctionfull = 
\forn{i}{\numberpredictionfunctions} c_i \predictionfunction_i(x),
\end{equation*}
\begin{equation*}
\forn{i}{\numberpredictionfunctions} c_i = 1; \quad c_i \geq 0, \quad i = \many{}{}{\numberpredictionfunctions};
\end{equation*}
выполнено
\begin{equation}\label{OptimizedDifferenceDecomposition}
\expectation \left[\left( \predictionfunction(x) - y \right)^2 \right] =
\forn{i}{\numberpredictionfunctions} c_i
\expectation \left[ \left(
	\predictionfunction_i(x) - y
\right)^2 \right] 
- \dfrac{1}{2}
\forn{i}{\numberpredictionfunctions} \forn{j}{\numberpredictionfunctions} c_i c_j
\expectation \left[ \left(
	\predictionfunction_i(x) - \predictionfunction_j(x)
\right)^2 \right]
\end{equation}

Соответственно, для уменьшения среднеквадратичной ошибки необходимо уменьшать среднеквадратичную ошибку каждого предиктора $\predictionfunction_i(x)$, а также увеличивать расхождение между прогнозами различных предикторов.

\subsection{Существующие методы}

В большинстве случаев функция-ансамбль $\predictionfunction(x) = \ensemblefunctionfull$ является линейной комбинацией функций предикторов $\many{\predictionfunction_}{(x)}{\numberpredictionfunctions}$. Коэффициенты линейной комбинации ищутся при помощи методов регрессионного анализа:
\begin{itemize}
\item гребневая регрессия~\cite{Ridge}
\item метод Лассо~\cite{Lasso}
\item эластичная сеть~\cite{ElasticNet}
\item регрессионная модель, отбирающая признаки, наиболее коррелирующие с откликом~\cite{ConvexCombinationsBestCorrelatedWithResponse}
\end{itemize}

Функции $\many{\predictionfunction_}{(x)}{\numberpredictionfunctions}$ в таких случаях обычно независимы друг от друга и получены методом оптимизации некоторого функционала.

\section{Ансамбль промежуточных решений}

Пусть:\\
$\objects$ - матрица из $\numberobjects$ строк, $i$-ая строка равна $x_i$;\\
$\results$ - вектор из $\numberobjects$ элементов, $i$-ый элемент равен $y_i$.

Пусть $\predictionfunction_1(x)$ представима в виде функции с параметрами $\optimizationmethodfunction(x, \theta)$ и была получена некоторым методом оптимизации функционала среднеквадратичной ошибки
$\mse(\predictionfunction_1(x), \objects, \results) = \dfrac{1}{\numberobjects}\forn{i}{\numberobjects} (\predictionfunction_1(x_i) - y_i)^2$:
$$\predictionfunction_1(x) = \optimizationmethodfunction(x, \theta_1)$$

В качестве функции $\optimizationmethodfunction(x, \theta)$ может выступать нейронная сеть, тогда параметрами $\theta$ будут являться веса нейронной сети.

\subsection{Поиск промежуточных решений}

Будем искать функции $\predictionfunction_2(x), \predictionfunction_3(x), \dots, \predictionfunction_\numberpredictionfunctions(x)$ в виде $\optimizationmethodfunction(x, \theta)$, минимизируя следующий функционал:
\begin{equation}\label{equation:lossfunction}
\loss(\predictionfunction_k(x), \objects, \results) = \mse(\optimizationmethodfunction(x, \theta_k), \objects, \results) -
\dfrac{\distinguishparameter}{k - 1}
\forn{i}{k - 1}\norm{\theta_k - \theta_i}^2, \quad k = 2, 3, \dots, \numberpredictionfunctions,
\end{equation}
где $\distinguishparameter \geq 0$ является гиперпараметром.

Данный функционал поощряет функции $\predictionfunction_i(x)$ иметь различные параметры $\theta_i$. Таким образом, мы пытаемся добиться максимального расхождения значений функций. Впоследствии это уменьшит среднеквадратичную ошибку функции-ансамбля $\predictionfunction(x)$, что следует из уравнения \reference{OptimizedDifferenceDecomposition}.

Стоит отметить, что при больших $k$ вычисление данного функционала может быть вычислительно затратным, так как $\forn{i}{k - 1}\norm{\theta_k - \theta_i}^2$ в общем случае требует $O(k)$ времени. Если все $\theta_i$ лежат в некотором евклидовом пространстве, данное выражение переписывается в виде:
$$
\forn{i}{k - 1}\norm{\theta_k - \theta_i}^2 = (k - 1)\norm{\theta_k}^2 - 2\scalarproduct{\theta_k, \forn{i}{k - 1} \theta_i} + \forn{i}{k - 1}\norm{\theta_i}^2.
$$

Значения выражений $\forn{i}{k - 1} \theta_i$ и $\forn{i}{k - 1}\norm{\theta_i}^2$ можно поддерживать в течении всего процесса поиска функции $\predictionfunction_k(x)$ и пересчитывать за $O(1)$ времени при переходе к поиску~$\predictionfunction_{k+1}(x)$.

\subsection{Поиск функции-ансамбля}

Составим матрицу $\predictedobjects$ размера $\numberobjects \times \numberpredictionfunctions$, где $\predictedobjects_{i, j} = \predictionfunction_j(x_i)$.
Для поиска функции-ансамбля $\predictionfunction(x) = \ensemblefunctionfull$ необходимо найти функцию $\ensemblefunction(\many{\hat{x}_}{}{\numberpredictionfunctions}) = \ensemblefunction(\hat{x})$, которую будем искать
из минимизации функционала $\mse(\ensemblefunction(\hat{x}), \predictedobjects, \results)$.

\section{Вычислительные эксперименты}

Для проверки качества работы представленного метода решим задачу классификации на трёх выборках данных систолического артериального давления. Данные заранее разбиты на обучающую и тестовую выборки. Размеры каждой из выборок представлены в таблице \ref{tabular:datasize}. Все параметры являются численными.

\begin{table}[h!]
	\begin{center}
		\begin{tabular}{|c|c|c|c|}
			\hline
			\multirow{2}{*}{Номер выборки} & \multicolumn{2}{c|}{Количество объектов} & \multirow{2}{*}{Количество параметров} \\
			\cline{2-3}
			& Обучающая выборка & Тестовая выборка & \\
			\hline
			1 & 718 & 458 & 114 \\
			\hline
			2 & 373 & 221 & 114 \\
			\hline
			3 & 832 & 495 & 114 \\
			\hline
		\end{tabular}
	\end{center}
	
	\caption{Размеры выборок данных}
	\label{tabular:datasize}
\end{table}

В качестве функции $\optimizationmethodfunction(x, \theta)$ используется двухслойная нейронная сеть с сигмоидальной функцией активации, которая оптимизируется методом Adam. В качестве функции-ансамбля $\ensemblefunction(\hat{x})$ используется точно такая же нейронная сеть, с другим количеством входных параметров и другой размерностью скрытого слоя.

В ходе эксперимента изначально обучается одна нейронная сеть $\predictionfunction_1(x) = \optimizationmethodfunction(x, \theta_1)$ и оценивается качество её работы. После этого обучается $\numberpredictionfunctions$ новых нейронных сетей $\predictionfunction_i(x) = \optimizationmethodfunction_i(x, \theta_i)$, а также ищется оптимальная функция-ансамбль $\ensemblefunction(\hat{x})$ в виде двухслойной нейронной сети. Затем оценивается качество работы ансамбля промежуточных решений $\predictionfunction(x) = \ensemblefunctionfull$.

\subsection{Результаты эксперимента}

Наилучших результатов удалось достичь для $\numberpredictionfunctions = 4$. Также при построении ансамбля промежуточных решений был замечен эффект улучшения качества предсказания отдельных нейронных сетей, обученных после первой нейронной сети $\predictionfunction_1(x)$.
Итоговые результаты эксперимента представлены в таблице \ref{tabular:results}.

\begin{table}[h!]
	\begin{center}
		\begin{tabular}{|c|c|c|}
			\hline
			Номер выборки & Одна нейронная сеть & Ансамбль промежуточных решений \\
			\hline
			1 & 0.78 & 0.86 \\
			\hline
			2 & 0.79 & 0.83 \\
			\hline
			3 & 0.78 & 0.81 \\
			\hline
		\end{tabular}
	\end{center}
	
	\caption{Результаты экспериментов (ROC AUC на тестовой выборке)}
	\label{tabular:results}
\end{table}

\subsection{Анализ полученных результатов}

Полученные результаты показывают, что ансамбль промежуточных решений, построенных с использованием функции потерь \reference{equation:lossfunction}, способен достаточно серьёзно увеличивать обобщающую способность предсказания в сравнении с отдельными нейронными сетями.
Также замечен эффект улучшения качества отдельных нейронных сетей, построенных данным методом. Попробуем объяснить данный эффект. Предполагается, что он вызван тем, что в случае нахождения недостаточно минимизирующего параметра $\theta_1$ для нейронной сети $\predictionfunction_1(x, \theta_1)$, последующие нейронные сети находят параметры $\theta_i$ на отдалении от $\theta_1$. Следовательно, $\theta_i$ не находится в области $\theta_1$, в которой функция потерь не достигает своего минимального значения.

\section{Заключение}

В процессе выполнения работы были получены следующие результаты:
\begin{itemize}
\item Был разработан метод повышения эффективности обучения, основанный на ансамбле промежуточных решений.
\item Были проведены вычислительные эксперименты, которые показали возможную применимость данного метода для улучшения качества нейронных сетей, решающих задачу классификации на реальных данных.
\item В ходе выполнения эксперимента было замечено улучшения качества работы отдельных нейронных сетей.
\end{itemize}

\end{document}
